
\section{Contraste bibliográfico: recirculación local en Couette--Poiseuille}
\label{sec:ref26-couette}

En un canal de separación \(h\), pared inferior fija y superior que se mueve a velocidad \(U>0\), sea \(G=-dp/dx\) el \emph{gradiente de presión impulsor}. Para fluido newtoniano de viscosidad constante, régimen estacionario, bidimensional, completamente desarrollado y sin deslizamiento, el balance de momento \(0=-dp/dx+\mu\,d^2u/dy^2\) da
\begin{equation}
u(y)=U\frac yh+\frac{G}{2\mu}\,y(h-y),\quad
q=\int_0^h u(y)\,dy=\frac{Uh}{2}+\frac{Gh^3}{12\mu}.
\label{eq:ref26-couette-q}
\end{equation}
La velocidad puede cambiar de signo cerca de la pared inferior aunque el \emph{caudal neto} conserve el signo positivo. Definiendo \(P=Gh^2/(2\mu U)\), la pendiente en \(y=0\) es \(du/dy|_0=(U/h)(1+P)\). La inversión local exige \(P<-1\), mientras que el caudal neto es positivo siempre que \(P>-3\). Así, \(-3<P<-1\) describe exactamente un flujo neto hacia delante con contraflujo junto a la pared fija. Este resultado es algebraico y no demuestra, por sí solo, que el flujo real sea estable ante perturbaciones tridimensionales.

\begin{tcolorbox}[title={Ejemplo guiado original: caudal y esfuerzos}, studiumdef, breakable]
Con \(h=0{,}002\,\mathrm{m}\), \(U=0{,}4\,\mathrm{m\,s^{-1}}\), \(\mu=0{,}020\,\mathrm{Pa\,s}\) y \(G=-8000\,\mathrm{Pa\,m^{-1}}\), resulta \(P=-2\) y \(dp/dx=+8000\,\mathrm{Pa\,m^{-1}}\), un gradiente adverso. La distribución es
\[
 u(y)=200\,y(1000y-1)\quad\text{(en SI, con \(y\) en metros)}.
\]
En \(0<y<0{,}001\,\mathrm{m}\) el fluido retrocede; en la otra mitad avanza. Sin embargo,
\[
 q=\frac{0{,}4\cdot0{,}002}{2}
    -\frac{8000(0{,}002)^3}{12\cdot0{,}020}
   =\mathbf{1{,}3333\cdot10^{-4}\,\mathrm{m^2\,s^{-1}}}>0.
\]
Los esfuerzos \emph{internos} \(\tau_{xy}=\mu\,du/dy\) en ambas paredes son \(\tau_{xy}(0)=-4\,\mathrm{Pa}\) y \(\tau_{xy}(h)=+12\,\mathrm{Pa}\). Al calcular fuerzas \emph{sobre las paredes}, debe utilizarse la normal y el signo correspondientes: no son automáticamente esos mismos vectores de tracción. La comprobación dimensional de \(q\) arroja \(\mathrm{m^2\,s^{-1}}\), al tratarse de caudal por unidad de anchura.
\end{tcolorbox}

\textbf{Autocomprobación.} Obtenga los valores de \(P\) que producen (i) pendiente nula en la pared fija, (ii) caudal neto nulo y (iii) ausencia de gradiente de presión. Justifique cuándo la condición de no deslizamiento y el supuesto de flujo desarrollado podrían fallar. Deducción y datos numéricos originales, contrastados conceptualmente con White, secciones 4.10 y 6,4 \cite{ref26-white}, y Shaughnessy, Katz y Schaffer \cite{ref26-shaughnessy}.
